\chapter{ソフトウェアによる MLP の実装}

\section{はじめに}

【目的】

本章では MNIST データセットを用いて訓練した MLP による手書き文字認識モデルをソフトウェアで構築します。
なお、本章は2名で実施することを想定しています。
適宜作業を分担してください。

【目標】

\begin{itemize}
    \item MNIST データセットを用いて MLP を訓練し、ソフトウェアによる識別モデルを作成できること。
    \item MNIST データセットを縮減し、訓練後の重み付けから抵抗値へ変換することで、効率のよいハードウェア実装に求められる要件について理解すること。
    \item ハードウェア駆動用の治具の作成を通して、組込み機器でのソフトウェア制御について理解すること。
\end{itemize}

【章構成】

\begin{enumerate}
    \item MNIST データセットを用いた MLP の訓練
    \item 訓練済みモデルのパラメタから抵抗値への変換
    \item Raspberry Pico を用いたテストベクタ出力機の作成
\end{enumerate}